\section{A common component in each one-particle slice}
\label{geom:section}

We now construct the retained sets required by
Lemma~\ref{setup:phase-transfer}. With the bath fixed, we cover its
forbidden balls by microscopic cubes and fill the bounded pockets enclosed
by each bad cluster. The retained points will lie in one component of the
open allowed slice. To bound the removed volume, we separate clusters
whose coordinate spans are at most one from larger clusters. The former
cost at most $Cr^2N$ in total filled volume. Each larger cluster forces order
$1/r$ distinct nearby bath particles with close partners, so the joint
spatial estimates control its occurrence. This argument requires no
connectivity of the full configuration space.

We retain the scaled torus $\Lambda_K$, its volume $V=K^3$, and the centered
unit cells $B_v$, where $v\in(\Z/K\Z)^3$. Write
\[
 Y=(x_2,\ldots,x_N),\qquad
 m(Y)=\int_{\Lambda_K}\Phi(x,Y)^2\dd x,\qquad
 \nu(\dd Y)=m(Y)\dd Y.
\]
Thus $\nu$ is the bath marginal of $\mu(\dd X)=\Phi(X)^2\dd X$.
Call a bath admissible if $d(x_i,x_j)>r$ for $2\le i<j\le N$.
Other baths have $m(Y)=0$. For an admissible bath, its open
one-coordinate slice is
\begin{equation}\label{geom:allowed-slice}
 S_Y=\Lambda_K\setminus\bigcup_{j=2}^N\overline B(x_j,r).
\end{equation}
We first work deterministically with admissible baths. The probability
estimate at the end uses the joint close-neighbor bound from
Equation~\eqref{sp:close-tail}.

\subsection{Microcubes and filled components}

Put
\[
 q_r=\lceil1/r\rceil,\qquad h=q_r^{-1},\qquad J=Kq_r.
\]
For $r\le1$, we have $r/2\le h\le r$. Subdivide every unit cell into
$q_r^3$ microcubes. In the periodic lift, their closed cubes are
\[
 Q_z=h\bigl(z+[0,1]^3\bigr)-(1/2,1/2,1/2),\qquad z\in\Z^3.
\]
Their index period is $J$ in every coordinate. A microcube is \emph{bad}
if its closure intersects a closed ball in~\eqref{geom:allowed-slice};
otherwise it is \emph{good}. Marking is periodic. Distinct index sites are
$*$-adjacent when their maximum coordinate distance is one. For a finite
bad $*$-component $C$, let $H_\infty(C)$ be the unique infinite
nearest-neighbor component of $\Z^3\setminus C$, and put
\[
 \operatorname{Fill}(C)=\Z^3\setminus H_\infty(C).
\]
Thus the fill adds all finite nearest-neighbor holes to $C$.
Define $\cA(Y)\subset\Lambda_K$ as the projection of all closed cubes with
indices in the fills of the bad components. If the lifted bad set has an
infinite component, define $\cA(Y)=\Lambda_K$. Use this convention also
for inadmissible baths. The retained region is
\[
 G(Y)=\Lambda_K\setminus\cA(Y).
\]

The fills remove bounded pockets behind the obstacles. This is the same
discrete filling operation used in Lemma~\ref{lat:boundary}; here its scale
is comparable to the exclusion distance. Phase transfer requires the
points of $G(Y)$ to lie in one component of $S_Y$; a connecting path
may use any good microcubes. Figure~\ref{geom:figure} distinguishes bad
cubes, their fill, and such a route.

\begin{figure}[t]
\centering
\begin{tikzpicture}[scale=.51]
 \foreach \x/\y in {3/2,4/2,5/2,6/2,3/3,6/3,3/4,6/4,3/5,4/5,5/5,6/5}
   \fill[black!32] (\x,\y) rectangle +(1,1);
 \fill[orange!28] (4,3) rectangle (6,5);
 \draw[black!20,step=1] (0,0) grid (10,8);
 \draw[very thick,blue!65!black,-{Stealth[length=5pt]}]
   (1.5,3.5)--(2.5,3.5)--(2.5,6.5)--(7.5,6.5)--(7.5,3.5)--(8.5,3.5);
 \fill[blue!65!black] (1.5,3.5) circle (.12) (8.5,3.5) circle (.12);
 \node[below left] at (1.5,3.5) {$x$};
 \node[below right] at (8.5,3.5) {$y$};
 \node[anchor=west] at (11,6) {\small bad microcubes};
 \fill[black!32] (10.25,5.7) rectangle +(0.5,0.5);
 \node[anchor=west] at (11,4.5) {\small filled pocket};
 \fill[orange!28] (10.25,4.2) rectangle +(0.5,0.5);
 \node[anchor=west] at (11,3) {\small route in good cubes};
 \draw[very thick,blue!65!black] (10.2,3)--(10.8,3);
\end{tikzpicture}
\caption{A planar schematic of the filling operation. The argument uses
three-dimensional microcubes. Each closed good cube avoids every closed
forbidden ball, so a finite route through good cubes has positive clearance.
The picture concerns a single-coordinate slice with the bath fixed.}
\label{geom:figure}
\end{figure}

\begin{lemma}[Connectivity with strict clearance]\label{geom:connectivity}
For every bath $Y$, any two points of $G(Y)$ can be joined by a continuous
path in $S_Y$ whose distance from the forbidden set is strictly positive.
\end{lemma}

\begin{proof}
There is nothing to prove when $G(Y)$ is empty. Otherwise the bath is
admissible and all lifted bad components are finite. Choose lifted microcube
sites containing the two points. Each lies outside the fill of every lifted bad component, and hence
in its infinite nearest-neighbor complement component.

Start with any finite nearest-neighbor path between these sites. If it
visits a bad component $C$, let $u$ precede the first visit and let $w$
follow the last visit. The path prefix before $u$ avoids $C$ and connects
$u$ to an endpoint outside its fill. The reversed suffix does the same for
$w$. Thus both vertices belong to the exterior $*$-boundary of $C$ visible
from infinity by nearest-neighbor paths. Lemma~\ref{lat:boundary} joins
them by a nearest-neighbor path in this boundary. All its vertices are
good: a bad $*$-neighbor of $C$ would already be in $C$.

This replacement removes bad vertices and introduces none. Repetition
therefore terminates, since the original path had finitely many bad
vertices. First and last visits are taken for one fixed lifted component;
different periodic copies are treated separately. Projecting the resulting
good-cube path gives a route on the torus.

Join successive cube centers through their common faces, and join the
original points to their cube centers. A point on a grid face can be
assigned any adjacent cube: membership outside $\cA(Y)$ ensures that each
such cube is good and outside every fill. Every closed good cube is
disjoint from every closed forbidden ball. The finite union of cubes used
by the route is compact, so its distance from the union of the closed
forbidden balls is positive. All the line segments stay in these cubes. This proves both
open-slice connectivity and strict clearance.
\end{proof}

\subsection{Small fills and distinct close-particle witnesses}

Call a finite microcomponent $C$ \emph{small} if
$h\diam_\infty C\le1$. All other components, including infinite ones,
are called \emph{large}. For a small fill, each occupied column has
cross-sectional area $h^2$ and length at most order one; charging each
occupied column to a bad cube gives the deterministic volume bound below. Large
components will instead force the close-neighbor events already estimated
under $\mu$.

\begin{lemma}[Volume of small fills]\label{geom:small-fills}
For $K>10$ and $r\le1$, the projected union of small fills has volume at
most $C r^2 N$, with an absolute constant $C$.
\end{lemma}

\begin{proof}
The fill of a finite component lies in its coordinate bounding box by
Lemma~\ref{lat:boundary}. Moreover, every column parallel to the first
coordinate axis that meets the fill also meets the component. Otherwise
that entire column would avoid the component and connect each of its
vertices to infinity, precluding a finite hole there.

For a small component $C$, there are at most $|C|$ such columns. Each
contributes length at most $1+2h$ in physical units and cross-sectional
area $h^2$. Hence its filled cubes have volume at most
\begin{equation}\label{geom:column-volume}
 (1+2h)h^2|C|\le3h^2|C|.
\end{equation}
A closed radius-$r$ ball meets at most $C_0=343$ microcubes, since
$r/h\le2$ and in each coordinate at most seven closed grid intervals
can meet its projection.

Choose one representative of each small component modulo the torus
period. Its projection is injective because each physical coordinate span
is at most one and $K>10$. The projected components are disjoint, and their
total number of bad microcubes is at most $C_0(N-1)$. Summing
Equation~\eqref{geom:column-volume}, and allowing overlaps of their
fills in this upper bound, gives $3C_0h^2(N-1)\le C r^2N$.
\end{proof}

\begin{lemma}[Distinct witnesses in a large component]
\label{geom:large-witnesses}
There are absolute constants $r_1,c_{\rm g}>0$ such that the following
holds for $r\le r_1$ and $K>10$. If a unit cell $B_v$ contains a
microcube from a large bad component, then at least $c_{\rm g}/r$
distinct bath particles lie within distance $5$ of $v$, and each has a
distinct bath particle within distance $6r$.
One may take $c_{\rm g}=1/(8C_0)$ with $C_0=343$.
\end{lemma}

\begin{proof}
Fix a lifted microcube index $z$ in that component and in a lift of $B_v$.
There is a component vertex at scaled maximum distance greater than
$1/2$ from $z$. For a finite component this follows from its diameter
being greater than one; for an infinite component it follows from local
finiteness. Follow a simple $*$-path from $z$, stopping when its scaled
maximum distance first reaches $1/4$. It contains at least $1/(4h)$
distinct vertices and remains within scaled distance $1/4+h$ of $z$.

Choose one bath ball witnessing badness at each path vertex. A fixed ball
witnesses at most $C_0$ of these vertices. All the chosen centers are in
a fixed small neighborhood of $v$, contained in its radius-$5$ ball for
sufficiently small $r$. This neighborhood and the cubes under consideration
have diameter less than $K$, so no two periodic copies of the same label
occur. The witness sequence therefore contains at least
$1/(4C_0h)\ge1/(4C_0r)$ distinct labels. Shrink $r_1$ so this number is
at least two, and retain the stated smaller constant $c_{\rm g}$.

Repeated labels do not destroy the close-partner conclusion. Form a graph
whose vertices are the distinct witness labels, joining different labels
whenever they occur consecutively in the sequence. This graph is connected
and has at least two vertices, so none is isolated. Two adjacent closed
microcubes meet. The centers of balls intersecting those cubes are therefore
at distance at most
\[
 2r+2\sqrt3\,h\le6r.
\]
Every distinct label thus has a different close partner, including labels
that recur in separated portions of the witness sequence.
\end{proof}

We now fix the geometric restriction on the close-count threshold:
\begin{equation}\label{geom:theta-choice}
 0<\theta<\frac{c_{\rm g}}{2\alpha_0}.
\end{equation}
This restriction is imposed before the choice of $b$ in the spatial
estimates. Since $r=\alpha/D$ and $\alpha\le2\alpha_0$, increasing
the lower threshold on $D$ ensures
\begin{equation}\label{geom:scale-separation}
 6r<s_0=bD^{-1/3},\qquad r\le r_1.
\end{equation}
For a full configuration $X=(x,Y)$, recall the close-count set from
Equation~\eqref{sp:close-bad-set}:
\begin{equation}\label{geom:coarse-bad-set}
 \mathcal Q(X)
 =\left\{v:\#\left\{i:d(x_i,v)\le5,
                    \min_{j\ne i}d(x_i,x_j)\le s_0\right\}
                   \ge\theta D\right\}.
\end{equation}
By Equation~\eqref{sp:close-tail}, for every deterministic set $S$ of
distinct coarse sites,
\begin{equation}\label{geom:coarse-joint-bound}
 \mu\{S\subset\mathcal Q(X)\}\le p_D^{|S|},
 \qquad p_D\longrightarrow0.
\end{equation}
Lemma~\ref{geom:large-witnesses} and
Equations~\eqref{geom:theta-choice}--\eqref{geom:scale-separation}
show that every parent cell of a large microcomponent belongs to
$\mathcal Q(X)$. Indeed $c_{\rm g}/r>\theta D$, and adding
the distinguished point $x$ to the bath can only decrease the relevant
nearest-neighbor distances. This implication holds for every $x$; it does
not identify $\nu$ with an $(N-1)$-particle ground-state law.

\subsection{Periodic animals and the omitted volume}

\begin{proposition}[Small omitted volume]\label{geom:omitted-volume}
Fix the constants as above and take $D$ sufficiently large. There are
absolute constants $C,C_{\rm an}$ such that, for $K\ge100$ and
$C_{\rm an}p_D\le1/2$,
\begin{equation}\label{geom:volume-bound}
 \delta_{\rm geom}:=\int\frac{|\cA(Y)|}{V}\,\nu(\dd Y)
 \le C r^2D+C p_D
       +K^3(C_{\rm an}p_D)^{\lfloor K/10\rfloor}.
\end{equation}
The retained points belong to one component of the open slice
$S_Y$ whenever the retained set is nonempty. The estimate is uniform
over the normalized ground vector and $\alpha\in[\alpha_0/2,2\alpha_0]$.
\end{proposition}

\begin{proof}
We estimate the bath-dependent volume by integrating under the full
configuration law:
\begin{equation}\label{geom:bath-volume-identity}
 \delta_{\rm geom}
 =\int_{\Lambda_K^N}\frac{|\cA(Y)|}{V}\,\mu(\dd x\,\dd Y).
\end{equation}
This is the marginal identity $\int\Phi(x,Y)^2\dd x=m(Y)$.
The coarse bad events used next may therefore depend on $x$ as well
as on the bath.

Set $m_K=\lfloor K/10\rfloor$ and let $\mathcal E_K$ be the event
that the coarse bad set in~\eqref{geom:coarse-bad-set} has a
$*$-component of size at least $m_K$. In a graph of degree $26$, the
number of connected vertex sets of size $m$ containing a prescribed root
is at most $C_{\rm an}^m$. To see this, choose for each set a deterministic
spanning tree and its depth-first traversal: the resulting walk has
$2(m-1)$ steps, and its visited set recovers the original set. There are
at most $26^{2(m-1)}$ such walks. A connected set of size at least $m_K$
contains a connected subset of size exactly $m_K$, by taking the first
$m_K$ vertices in a tree exploration. Equation~\eqref{geom:coarse-joint-bound}
and a union bound over roots give
\begin{equation}\label{geom:large-animal-error}
 \mu(\mathcal E_K)
 \le K^3(C_{\rm an}p_D)^{m_K}.
\end{equation}

On $\mathcal E_K^c$, every coarse bad component has fewer than $K/10$
vertices. We make its periodic lifting explicit. Lift a spanning tree of
a component with $m$ vertices. Each additional edge and its tree path
form a loop of length at most $2m-1<K$. Such a loop has zero winding:
every step changes each coordinate by at most one, whereas nonzero winding
requires displacement of magnitude at least $K$ in one coordinate.
Every additional edge therefore closes in the chosen finite tree lift.
This proves that the full lifted component is finite, projects bijectively,
and has coordinate span at most $m-1$. It is the nearest-range case of
Lemma~\ref{lat:small-lifts}.

The parent-cell map sends adjacent microcubes to equal or $*$-adjacent
unit cells. Every large microcomponent is therefore contained in the unit
cubes of one lifted coarse bad component. In particular it is finite on
$\mathcal E_K^c$: each of those finitely many unit cubes contains only
$q_r^3$ microcubes. Its fill lies in its own coordinate bounding box and
thus in the coordinate box of those unit cubes.

Consequently, if a prescribed unit cell is affected by a large fill,
there is a coarse bad animal of some size $m\ge1$ with a root within
maximum distance $C m$ of that cell. There are at most $C m^3$ possible
roots, even after projecting to the torus. The rooted-animal count and
Equation~\eqref{geom:coarse-joint-bound} give
\begin{equation}\label{geom:filled-animal-sum}
 \mu\{\mathcal E_K^c\text{ and the cell meets a large fill}\}
 \le\sum_{m\ge1} C m^3(C_{\rm an}p_D)^m\le C'p_D.
\end{equation}
The last inequality follows by summing the convergent power series on
$C_{\rm an}p_D\le1/2$. No independence of coarse sites is used.

Sum Equation~\eqref{geom:filled-animal-sum} over unit cells and divide
by $V$. On $\mathcal E_K$, bound the discarded fraction by one.
Lemma~\ref{geom:small-fills} bounds the small-fill contribution by
$Cr^2N/V=Cr^2D$. Combining these bounds with
Equation~\eqref{geom:large-animal-error} proves
Equation~\eqref{geom:volume-bound}, by
Equation~\eqref{geom:bath-volume-identity}.
Lemma~\ref{geom:connectivity} gives the connectivity assertion.

All sets used here are measurable. Bad-cube membership is a finite family
of distance tests in one period. A finite periodic component has a finite
lift precisely when its spanning-tree lift closes along every additional
edge. In that case its fill is determined by connectivity in a finite
bounding box, with the exterior attached. These finite procedures determine
the projected discarded cubes and their volume from the bath positions.
\end{proof}

The right side of Equation~\eqref{geom:volume-bound} has the needed order
of limits. Since $r^2D=\alpha^2/D$, its first two terms can be made
arbitrarily small by increasing the lower threshold on $D$. At each such
fixed density, the last term tends to zero as $K$ grows. Its displayed
bound also gives uniform convergence for all $D$ above a threshold at
which $C_{\rm an}p_D\le1/2$. The estimate is therefore ready for the
phase-transfer lemma and the final parameter choice.
